%----------------------------------------------------------------------
\chapter{Introduction}
\label{ch:intro}
%----------------------------------------------------------------------

You may have read about similar things in \cite{Goodliffe2007}.
You can also write footnotes.\footnote{Footnotes will be positioned automatically.}
You can reference acronyms. \gls{oci} prints it in longform, the second reference \gls{oci} is only the acronym.

\blindtext

\blindtext

%----------------------------------------------------------------------
\section{This is an Important Section}
\label{sec:section}
%----------------------------------------------------------------------

It is possible to reference glossary entries as \gls{library} and \gls{ast} as an example.

\blindtext

%----------------------------------------------------------------------
\subsection{And an even more important subsection}
\label{sec:subsection}
%----------------------------------------------------------------------

\blindtext
